% Part: first-order-logic
% Chapter: beyond
% Section: modal-logics

\documentclass[../../../include/open-logic-section]{subfiles}

\begin{document}

\olfileid{fol}{byd}{mod}

\olsection{మోడల్ తర్కాలు}

కింది సోపాధిక వాక్య ఉదాహరణను పరిశీలించండి:
\begin{quote}
  జెరెమీ ఆ గదిలో ఒంటరిగా ఉంటే, అతను మద్యం మత్తులో, నగ్నంగా,
  కుర్చీలపై నృత్యం చేస్తున్నాడు.
\end{quote}
ఇది ద్రవ్యపరంగా సత్యమై ఉండగలిగినా తప్పుదారి పట్టించే సోపాధిక
ప్రకటనకు ఉదాహరణ. ఎందుకంటే, పూర్వాంగం అసత్యం లేదా ఉత్తరాంగం సత్యం
అనేదాని కంటే పూర్వాంగానికీ ఉత్తరాంగానికీ మధ్య మరింత బలమైన సంబంధం
ఉందని ఇది సూచిస్తున్నట్లు అనిపిస్తుంది. అంటే, ఈ నిర్దిష్ట లోకంలోనే
ఆ ప్రకటన సత్యమని (జెరెమీ గదిలో ఒంటరిగా లేడు కాబట్టి ఇక్కడ అది
తుచ్ఛంగా సత్యం కావచ్చు) కాకుండా, పూర్వాంగం సత్యమై \emph{ఉంటే}
ఉత్తరాంగం కూడా సత్యమై \emph{ఉండేదని} ఆ మాటలు సూచిస్తాయి. మరో
మాటలో, ఈ లోకంలో మాత్రమే కాక ఏ ``సాధ్యమైన'' లోకంలోనైనా ప్రకటన
సత్యమని; లేదా ఈ నిర్దిష్ట లోకంలో కేవలం సత్యమవడానికి భిన్నంగా అది
\emph{అవశ్యంగా} సత్యమని ఆ ప్రకటనకు అర్థం చెప్పవచ్చు.

ఇటువంటి అవశ్యకతకు అర్థం చెప్పడానికి మోడల్ తర్కాన్ని రూపొందించారు.
సాధారణ ప్రతిజ్ఞావాక్య తర్కానికి ఒక పెట్టె సంయోజకాన్ని చేరిస్తే మోడల్
ప్రతిజ్ఞావాక్య తర్కం లభిస్తుంది; అంటే $!A$ \tetoken{ఒక సూత్రం}{!!a{formula}} అయితే,
$\Box !A$ కూడా ఒక సూత్రమే. అంతఃప్రజ్ఞాత్మకంగా, $!A$ \emph{అవశ్యంగా}
సత్యమని, లేదా సాధ్యమైన ప్రతి లోకంలో సత్యమని $\Box !A$ ప్రకటిస్తుంది.
$\Diamond !A$ను సాధారణంగా $\lnot \Box \lnot !A$కు సంక్షిప్త రూపంగా
తీసుకొని, $!A$ \emph{సాధ్యంగా} సత్యమని ప్రకటిస్తున్నట్లు చదువుతారు.
మోడాలిటీని విధేయ తర్కానికి కూడా చేర్చవచ్చు.

మోడల్ తర్కానికి అర్థవిచారం ఇవ్వడానికి క్రిప్కె \tetoken{నిర్మాణాలను}{!!{structure}s}
వాడవచ్చు; నిజానికి క్రిప్కె మొదట ఈ అర్థవిచారాన్ని మోడల్ తర్కాన్ని
దృష్టిలో పెట్టుకొని రూపొందించాడు. పాక్షిక క్రమాలకే పరిమితం కాకుండా,
మరింత సాధారణంగా ``సాధ్యమైన లోకాల'' సమితి~$P$, లోకాల మధ్య
ద్విస్థాన ``ప్రాప్యత'' సంబంధం $\Atom{R}{x,y}$ ఉంటాయి.
అంతఃప్రజ్ఞాత్మకంగా, లోకం~$q$ లోకం~$p$కు అనుకూలమని
$\Atom{R}{p,q}$ ప్రకటిస్తుంది; అంటే, మనం లోకం~$p$లో ``ఉంటే'', లోకం
$q$లా ఉండి ఉండే సాధ్యతను కూడా పరిగణించాలి.

విస్తారాత్మక (ఎక్స్టెన్షనల్) తర్కానికి భిన్నంగా మోడల్ తర్కాన్ని
కొన్నిసార్లు ``అంతర్భావాత్మక'' (ఇంటెన్షనల్) తర్కం అంటారు.
సాంప్రదాయిక తర్కం వంటి విస్తారాత్మక తర్కానికి ఉద్దేశిత అర్థవిచారం
``వాస్తవమైన'' ఒకే లోకాన్ని సూచిస్తుంది; అంతర్భావాత్మక తర్కపు
అర్థవిచారం మరింత విస్తృతమైన సత్తావిచారంపై ఆధారపడుతుంది. అవశ్యకతకు
\tetoken{నిర్మాణం}{!!{structure}} ఇవ్వడమే కాకుండా, అనువర్తనానికి అనుగుణంగా $\Box$,
$\Diamond$లను పునరర్థనిర్దేశం చేస్తూ ఇతర భాషా నిర్మితులకు
\tetoken{నిర్మాణం}{!!{structure}} ఇవ్వడానికి కూడా మోడాలిటీని వాడవచ్చు. ఉదాహరణకు:
\begin{enumerate}
\item నిరూప్యతా తర్కంలో $\Box !A$ను ``$!A$ నిరూపించదగినది'' అని,
  $\Diamond !A$ను ``$!A$ అవైరుధ్యమైనది'' అని చదువుతారు.
\item జ్ఞానమీమాంసాత్మక తర్కంలో $\Box !A$ను ``నాకు $!A$ తెలుసు'' లేదా
  ``నేను $!A$ను నమ్ముతున్నాను'' అని చదవవచ్చు.
\item కాలిక తర్కంలో $\Box !A$ను ``$!A$ ఎల్లప్పుడూ సత్యం'' అని,
  $\Diamond !A$ను ``$!A$ కొన్నిసార్లు సత్యం'' అని చదవవచ్చు.
\end{enumerate}

మోడాలిటీతో వ్యవహరించే నియమాలు, స్వీకృతాలను తర్కానికి చేర్చాలనుకుంటాం.
ఉదాహరణకు, $\Log{S4}$ వ్యవస్థలో ప్రతిజ్ఞావాక్య తర్కపు సాధారణ
స్వీకృతాలు, నియమాలతోపాటు కింది స్వీకృతాలు ఉంటాయి:
\begin{align*}
& \Box (!A \lif !B) \lif (\Box !A \lif \Box !B)\\
& \Box !A \lif !A \\
& \Box !A \lif \Box \Box !A
\intertext{అలాగే ``$!A$ నుంచి $\Box !A$ను నిగమించు'' అనే నియమం
  ఉంటుంది. $\Log{S5}$ కింది స్వీకృతాన్ని చేర్చుతుంది:}
& \Diamond !A \lif \Box \Diamond !A
\end{align*}
ఈ స్వీకృతాల భేదాలు వేర్వేరు అనువర్తనాలకు తగినవి కావచ్చు; ఉదాహరణకు,
తార్కిక అవశ్యకత భావనను S5 లక్షణీకరిస్తుందని సాధారణంగా
తీసుకుంటారు. క్రిప్కె \tetoken{నిర్మాణాలలోని}{!!{structure}s} ప్రాప్యత సంబంధాన్ని సహజమైన
విధాల్లో పరిమితం చేయడం ద్వారా, \tetoken{వ్యుత్పత్తి}{!!{derivation}} వ్యవస్థ యథార్థంగానూ
సంపూర్ణంగానూ ఉండే అర్థవిచారాన్ని సాధారణంగా కనుగొనగలగడం దీని మంచి
లక్షణం. ఉదాహరణకు, ప్రాప్యత సంబంధం పరావర్తనమూ సంక్రమణమూ అయ్యే
క్రిప్కె \tetoken{నిర్మాణాలు}{!!{structure}s} తరగతికి $\Log{S4}$ అనుగుణంగా ఉంటుంది.
ప్రాప్యత సంబంధం {\em సార్వత్రికం} అయ్యే క్రిప్కె \tetoken{నిర్మాణాలు}{!!{structure}s}
తరగతికి $\Log{S5}$ అనుగుణంగా ఉంటుంది; అంటే ప్రతి లోకం నుంచి ప్రతి
ఇతర లోకానికీ ప్రాప్యత ఉంటుంది. కాబట్టి ప్రతి లోకంలో $!A$ సత్యమైతేనే,
అప్పుడు మాత్రమే $\Box !A$ సత్యమవుతుంది.

\end{document}
